%---------------------------------------------------------------------
% Key names and dates, then the recurring mistakes.
%---------------------------------------------------------------------

\section*{اہم نام، اخبارات اور سنین}
\markboth{اہم نام اور سنین}{}

تاریخ والے سوال میں نمبر \emph{ناموں اور سنین} سے ملتے ہیں۔ یہ فہرست یاد کر
لیجیے --- سوال \textenglish{2} اور \textenglish{3} اسی پر کھڑے ہیں۔

\medskip
\begin{center}
\begin{tabular}{@{}p{4.2cm} p{2.4cm} p{7.8cm}@{}}
\toprule
\textbf{اخبار / رسالہ} & \textbf{سنہ} & \textbf{شخصیت اور اہمیت} \\
\midrule
\textenglish{Bengal Gazette} & \textenglish{1780} & \textenglish{J. A. Hicky} --- برصغیر کا پہلا اخبار \\
جامِ جہاں نما & \textenglish{1822} & پہلا اردو اخبار، کلکتہ \\
دہلی اردو اخبار & \textenglish{1837} & مولوی محمد باقر --- پہلے شہید صحافی \\
سائنٹفک سوسائٹی گزٹ & \textenglish{1866} & سر سید احمد خان --- جدید تعلیم \\
تہذیب الاخلاق & \textenglish{1870} & سر سید --- سماجی و فکری اصلاح \\
زمیندار & --- & مولانا ظفر علی خان --- بابائے صحافت \\
کامریڈ & \textenglish{1911} & مولانا محمد علی جوہر (انگریزی) \\
الہلال & \textenglish{1912} & مولانا ابوالکلام آزاد \\
ہمدرد & \textenglish{1913} & مولانا محمد علی جوہر (اردو) \\
البلاغ & \textenglish{1915} & مولانا ابوالکلام آزاد \\
انقلاب & \textenglish{1927} & لاہور \\
جنگ & \textenglish{1939} & میر خلیل الرحمٰن \\
نوائے وقت & \textenglish{1940} & حمید نظامی \\
ڈان & \textenglish{1941} & قائدِ اعظم کی رہنمائی میں \\
\bottomrule
\end{tabular}
\end{center}

\medskip
\begin{nukta}[چار اصطلاحات جو اکثر گڈ مڈ ہوتی ہیں]
\begin{center}
\begin{tabular}{@{}p{2.6cm} p{11.6cm}@{}}
\toprule
\textbf{خبر} & صرف واقعہ --- \emph{کیا، کب، کہاں، کون}۔ رائے شامل نہیں۔ \\
\textbf{اداریہ} & \emph{ادارے} کی رائے، بے نام، تازہ مسئلے پر، مختصر۔ \\
\textbf{کالم} & \emph{فرد} کی رائے، نام و تصویر کے ساتھ، باقاعدہ۔ \\
\textbf{فیچر} & تفصیلی، ادبی، انسانی دلچسپی کی تحریر؛ وقت کی پابند نہیں۔ \\
\bottomrule
\end{tabular}
\end{center}
\end{nukta}

\clearpage

%---------------------------------------------------------------------
\section*{آٹھ عام غلطیاں}
\markboth{آٹھ عام غلطیاں}{}

\begin{enumerate}\setlength{\itemsep}{5pt}
  \item \textbf{مسلسل پیراگراف میں جواب لکھنا۔} بیس نمبر کے جواب میں
        \emph{عنوانات} لگائیے اور نکات شمار کیجیے۔ ممتحن کو ڈھانچہ نظر آنا
        چاہیے۔
  \item \textbf{مثالیں نہ دینا۔} ہر جواب میں پاکستانی اخبارات، شخصیات یا
        واقعات کی مثال دیجیے۔ بغیر مثال کے جواب نصف رہ جاتا ہے۔
  \item \textbf{سنین یاد نہ ہونا۔} تاریخ کے سوال میں
        \textenglish{1780}، \textenglish{1822}، \textenglish{1866}،
        \textenglish{1912} جیسے سنین ہی امتیاز پیدا کرتے ہیں۔
  \item \textbf{اداریے اور کالم کو ایک سمجھنا۔} اداریہ ادارے کی بے نام رائے
        ہے، کالم فرد کی نام والی رائے۔
  \item \textbf{رپورٹر اور سب ایڈیٹر کے فرائض ملا دینا۔} رپورٹر خبر
        \emph{لاتا} ہے، سب ایڈیٹر اسے \emph{سنوارتا} ہے اور سرخی لگاتا ہے۔
  \item \textbf{قانون اور ضابطۂ اخلاق کو ایک سمجھنا۔} قانون ریاست نافذ کرتی
        ہے (\textenglish{PEMRA}، \textenglish{PECA})؛ ضابطۂ اخلاق صحافی خود پر
        لاگو کرتا ہے (پریس کونسل، \textenglish{APNS})۔
  \item \textbf{رائے والے سوال کا یکطرفہ جواب۔} ''کیا میڈیا آزاد ہے'' جیسے
        سوال میں پہلے دونوں پہلو، پھر اپنی رائے۔
  \item \textbf{وقت کی بدانتظامی۔} پانچ سوال، تین گھنٹے --- فی سوال تقریباً
        پینتیس منٹ۔ پہلے سوال پر ایک گھنٹہ لگا دینے سے آخری سوال ادھورا رہ
        جاتا ہے، اور یہی سب سے بڑا نمبر کھونے کا سبب ہے۔
\end{enumerate}

\vfill
\begin{center}
\begin{tcolorbox}[enhanced, colback=bmtint, colframe=bmblue, boxrule=0.5pt,
                  arc=3pt, width=0.94\linewidth,
                  left=11pt, right=11pt, top=9pt, bottom=9pt]
\small
\textbf{\color{bmblue}اگر وقت بہت کم ہو۔}\ \
دونوں پرچوں میں دہرائے گئے چار موضوعات سب سے پہلے تیار کیجیے:
\textenglish{(1)} صحافت کا ارتقاء (سوال \textenglish{2})؛
\textenglish{(2)} سب ایڈیٹر کے اوصاف و فرائض (سوال \textenglish{7})؛
\textenglish{(3)} ضابطۂ اخلاق اور صحافتی قوانین (سوال \textenglish{8}،
\textenglish{9})؛
\textenglish{(4)} اداریہ --- تعریف، اقسام اور فیچر سے فرق (سوال
\textenglish{4}، \textenglish{5})۔
پانچ میں سے تین سوال عموماً انہی سے بن جاتے ہیں۔

\medskip
\textbf{\color{bmblue}ماخذ۔}\ \
اس کورس کی سرکاری درسی کتاب دستیاب نہ ہو سکی، اس لیے تمام سوالات اے آئی او یو
کے دو حقیقی پرچوں سے لیے گئے ہیں: بی اے اصولِ صحافت (\textenglish{430})،
سمسٹر بہار \textenglish{2013} اور سمسٹر خزاں \textenglish{2013}۔ جوابات صحافت
کے مسلمہ اصولوں اور پاکستانی صحافت کی معلوم تاریخ کے مطابق ہیں --- کتاب کی
مخصوص تفصیلات اپنی درسی کتاب سے ملا لیجیے۔
\end{tcolorbox}
\end{center}
